% ==============================================================
% 第1章 绪论
% ==============================================================
\chapter{绪论}

\section{研究背景}

\subsection{中国人口老龄化与慢性病负担}

中国正经历全球规模最大、速度最快的人口老龄化进程。根据民政部2024年度国家老龄事业发展公报，截至2024年末，全国60周岁及以上老年人口已达31031万人，占总人口比例升至22.0\%；65周岁及以上老年人口22023万人，占比15.6\%。自2019年至2024年，60岁以上老年人口年均增长率高达3.7\%，远超全球平均水平。世界卫生组织预测，到2040年中国60岁以上人口规模将达4.02亿，届时每三个中国人中便有一位是老年人。随着人口结构深度转型，老年群体的健康管理压力急剧攀升。国家卫生健康委统计数据显示，我国超过1.9亿老年人患有慢性病，慢性病导致的死亡已占总死亡人数的88.5\%，且城市地区老年人多病共存现象尤为突出，平均每位老年人同时罹患2.3种慢性病\todocite{}\todocite{}\todocite{}。

在众多慢性病中，2型糖尿病（Type 2 Diabetes Mellitus，T2DM）是威胁老年人健康的重中之重。根据国际糖尿病联盟（IDF）2021年最新数据，中国20---79岁人群糖尿病患病率为10.6\%，患病人数达1.41亿，是全球糖尿病患者最多的国家。《中国糖尿病防治指南（2024版）》数据显示，我国糖尿病患病率自2013年的10.9\%上升至2018---2019年的12.4\%，各民族及各地区之间存在明显差异。尤为令人警惕的是，我国糖尿病长期处于``知晓率、治疗率、控制率''三低困境：知晓率仅36.7\%、治疗率仅32.9\%、控制率仅50.1\%，仍处于较低水平。滕卫平团队对7.59万名成年受试者的全国调查更发现，约56.7\%的糖尿病患者并不知道自己已患病，这意味着存在庞大的隐性患病人群，公共卫生防治压力极为严峻\todocite{}\todocite{}。

老年糖尿病患者还面临一系列独特的管理难题，使其区别于普通成人糖尿病患者。首先，低血糖风险大幅增加：研究表明，重度低血糖组患者的认知功能障碍发生率高达75.00\%，高频次低血糖（$\geq$8次）组认知功能障碍发生率也达68.97\%，低血糖的发生次数与程度是影响老年糖尿病患者认知水平的重要危险因素。其次，多病共存是常态，老年糖尿病患者往往合并高血压、冠心病、慢性肾病等多种疾病，药物间的相互作用和禁忌证使用药方案极为复杂。再者，认知功能下降会直接削弱患者的自我管理能力，导致服药依从性、血糖监测频率、饮食控制能力全面下降，形成``糖尿病---认知衰退---自我管理失效''的恶性循环。上述特殊性决定了老年糖尿病管理需要更个性化、更精细化、更具连续性的医疗支持体系\todocite{}\todocite{}\todocite{}。

\subsection{基层医疗体系面临的结构性困境}

面对庞大的老年慢病人群，国家层面积极推进分级诊疗体系建设，以期将慢性病管理重心下沉至基层。近年来国务院常务会议明确提出，要统筹推进分级诊疗体系建设和医疗卫生强基工程，以常见病、慢性病为重点引导群众基层首诊。截至2024年，全国基层医疗卫生机构数量已突破104万所，基层医务人员超525万名，超90\%的乡镇卫生院和社区卫生服务中心已达到基本服务能力标准。国家卫生健康委亦于2025年发布《关于促进和规范``人工智能+医疗卫生''应用发展的实施意见》，进一步推动医疗卫生领域大模型规范应用\todocite{}\todocite{}。

然而，政策推行与基层实际承接能力之间仍存在显著落差。以糖尿病管理为核心痛点来看，家庭医生签约服务在实践中普遍面临``签而不约、约而不管''的困境------服务覆盖面广但质量参差不齐，签约率指标与实质性健康管理之间存在明显鸿沟。系统性综述研究表明，家庭医生合同服务的有效性受到工作负荷过重、专业能力不足、随访机制不健全等多重障碍的制约。在具体管理指标上，深圳某社区卫生中心对359名2型糖尿病患者的随访研究显示，仅有62.70\%的患者达到血糖控制标准，相比临床期望存在明显差距。与此同时，院内诊疗数据与院外居家监测数据之间的``数字孤岛''问题突出：患者在社区随访中记录的血糖数据、用药情况与院内电子病历系统长期割裂，医生无法全面掌握患者的动态健康状况，导致随访流于形式、干预决策缺乏依据\todocite{}\todocite{}\todocite{}。

基层全科医生的能力瓶颈与工作负荷矛盾是上述困境的深层根源。基层医生不仅需要承担辖区内大量高血压、糖尿病、心脑血管疾病等慢病患者的定期随访任务，还要完成公共卫生服务、基本医疗诊疗等工作，繁重的工作量严重限制了对每位患者提供个性化、高质量管理服务的时间和精力。此外，基层医生在处理老年糖尿病患者的复杂并发症、多药联用风险以及制订个体化血糖控制目标等方面，专业能力相对不足，亟需有效的辅助决策工具支撑。这种``人力资源有限、管理需求无限''的结构性矛盾，是当前基层慢病管理陷入困境的核心逻辑\todocite{}。

\subsection{人工智能技术赋能基层医疗的政策机遇}

在政策层面，中国人工智能医疗的顶层设计日趋完善。2025年，国家卫生健康委联合多部门发布《关于促进和规范``人工智能+医疗卫生''应用发展的实施意见》，明确支持人工智能大模型研发和迭代升级，推动医疗卫生领域大模型规范备案。国家卫生健康委的战略愿景更为宏大：到2030年，AI辅助诊疗将在9亿中国公民的基层医疗机构中成为标准配置。据预测，中国AI医疗市场规模将从2020年的9亿美元增长至2023年的15.9亿美元，并有望在2030年前后达到188.8亿美元\todocite{}\todocite{}\todocite{}。

从技术发展轨迹来看，人工智能在医疗领域已完成从感知识别到认知推理的代际跃迁。早期医疗AI以图像识别为核心（如影像辅助诊断），而以大语言模型（Large Language Models，LLM）为代表的新一代AI具备处理非结构化文本、进行多步骤推理、理解临床语境的能力，开始向临床决策支持、患者教育、慢病管理等更广泛的认知型任务渗透。在此基础上，多智能体系统（Multi-Agent Systems，MAS）通过协调多个专业化Agent的分工协作，进一步突破了单一模型的能力边界，为复杂医疗场景中的多维决策提供了新的技术范式\todocite{}\todocite{}\todocite{}。

LLM与Agent技术在社区慢病管理领域的应用前景尤为令人期待。清华大学``Agent Hospital''的实践表明，由多个AI医生Agent构成的系统可实现对大量患者的高效诊疗服务。将上述技术路径应用于基层糖尿病管理，有望在以下维度产生实质性突破：通过RAG（检索增强生成）构建糖尿病专业知识库，使AI能够遵循最新临床指南提供准确建议；通过多Agent协同模拟``全科医生---专科顾问---患者管理''的分层决策结构；通过长时记忆机制为每位患者建立动态健康画像，实现真正的个性化管理。这些技术路径的有机整合，正是本研究的核心出发点\todocite{}。


% ------------------------------------------------------------------
\section{研究意义}

\subsection{理论意义}

本研究的理论价值体现在以下两个层面。

其一，\textbf{探索多智能体协同机制在垂直医疗场景中的设计范式}。现有多Agent框架（如MDAgents、RareAgents）主要聚焦于单次诊断决策任务，缺乏对慢病管理这一需要长期连续服务场景的系统性研究。本研究通过将多Agent架构应用于社区糖尿病长期随访管理，探索``全科Agent---专科顾问Agent---患者交互Agent---预警Agent''的协同决策机制，为LLM-based MAS在垂直医疗场景中的设计提供了可参考的架构范式，有助于填补MAS在慢病管理领域的研究空白\todocite{}\todocite{}。

其二，\textbf{为LLM驱动的慢病风险评估提供可参考的模型框架}。将提示词工程（Prompt Engineering）、检索增强生成（RAG）与个性化健康画像相结合，构建面向中国基层糖尿病管理实践的AI辅助风险评估框架，有助于推动医疗AI从通用问答向垂直领域精准决策支持的理论演进，同时为《中国糖尿病防治指南（2024版）》等临床规则的智能化转化提供方法论依据\todocite{}\todocite{}。

\subsection{实践意义}

本研究的实践价值体现在三个维度。

\textbf{提升社区糖尿病管理的效率与精细度}。通过构建多Agent智能随访系统，实现对患者血糖波动、并发症风险、用药依从性的自动化监测与分级预警，将基层医生从重复性随访记录工作中解放出来，使其能够将有限精力集中于高风险患者的精准干预。研究表明，标准化随访管理本身即可显著改善患者的血糖控制结局；而智能化系统通过提升随访质量与个性化程度，有望进一步突破当前血糖达标率的瓶颈\todocite{}。

\textbf{降低老年患者数字化健康管理的使用门槛}。老年糖尿病患者普遍面临技术焦虑、视力下降、操作不熟练等数字鸿沟障碍。本系统拟引入语音识别（ASR）交互模式，结合适老化界面设计，使老年患者可以通过自然语言完成血糖录入、健康咨询等操作，从根本上降低数字健康工具的使用门槛。可解释性AI（XAI）的设计思路亦有助于提升老年用户对AI推荐结果的信任度，促进健康管理行为的持续参与\todocite{}\todocite{}\todocite{}。

\textbf{辅助基层医生实现从``被动治疗''到``主动管理''的转变}。通过Human-in-the-loop模式，系统将AI生成的风险预警、随访建议、转诊提示以结构化形式呈现给基层医生，实现人机协同决策。这一模式不仅保障了医疗安全性，也能够在不显著增加工作负荷的前提下，帮助基层医生实现对辖区慢病患者的主动监测与提前干预，推动基层医疗服务模式从``坐等就诊''向``主动管理''转变\todocite{}\todocite{}。


% ------------------------------------------------------------------
\section{国内外研究现状}

\subsection{糖尿病社区管理模式研究现状}

\noindent\textbf{国内研究现状}

国内社区糖尿病管理主要依托家庭医生签约服务体系和慢病一体化门诊两种模式推进。家庭医生签约模式以居民主动签约为基础，由社区医生提供定期随访、血糖监测、用药指导等基本管理服务，被证明可有效改善糖尿病患者的就医连续性。然而，系统性综述表明，家庭医生合同服务的实际效果受到工作负荷过重、医生能力参差不齐、患者依从性低以及随访流于形式等多重障碍的制约，其管理效果在不同地区间存在显著差异。厦门市的实践探索形成了``专科医师---全科医生---健康管理师''三师共管的慢病分级管理模式，在提升管理质量方面取得一定成效，但依然面临人力资源不足与可复制性有限的挑战。此外，基于mHealth的干预研究表明，数字化健康工具有助于改善T2DM患者的危险因素控制率，在农村社区尤为明显，但这些系统仍高度依赖人工随访，个性化程度不足，难以形成可持续的全程管理闭环\todocite{}\todocite{}\todocite{}\todocite{}\todocite{}。

\noindent\textbf{国外经验借鉴}

英国的\textbf{质量与结果框架}（Quality and Outcomes Framework，QOF）为全科医生提供基于绩效的糖尿病管理财务激励机制，围绕HbA1c控制、血压管理、血脂监测等核心指标构建了系统化的质量评价体系，有力推动了英国基层糖尿病诊疗质量的持续提升。QOF的实践表明，将经济激励与标准化评估结合，可以有效驱动基层医生的主动管理行为，但也存在过度依赖指标导向、忽视个体化需求的内在张力。美国的\textbf{以患者为中心的医疗之家}（Patient-Centered Medical Home，PCMH）模型则以团队化照护、患者自我管理支持和数据驱动质量改进为核心，强调将糖尿病慢性病管理模型（Chronic Care Model）嵌入基层实践，多项研究证实PCMH能够改善HbA1c、血脂、血压等关键糖尿病管理指标\todocite{}\todocite{}\todocite{}\todocite{}\todocite{}\todocite{}。

\noindent\textbf{现有模式的共同瓶颈}

无论是中国的家庭医生签约模式，还是英国的QOF或美国的PCMH，均面临三方面共同瓶颈：第一，\textbf{高度人力依赖}------管理质量与医护人员的时间、能力高度绑定，难以规模化推广；第二，\textbf{个性化程度不足}------现有系统多以统一标准进行人群化管理，难以满足老年糖尿病患者高度异质化的个体需求；第三，\textbf{随访连续性差}------缺乏院内外数据打通机制，患者在院外的血糖变化、生活方式改变等动态信息无法及时传递给医生，导致随访干预缺乏针对性。这些共同瓶颈正是本研究希望通过LLM+MAS技术路径加以突破的核心问题域\todocite{}。

\subsection{医疗垂类大模型研究现状}

\noindent\textbf{通用LLM在医疗领域的局限性}

GPT-4等通用大语言模型在医学考试题目等静态基准测试中表现出色，但在实际临床应用中面临两大核心挑战：\textbf{幻觉（Hallucination）问题}和\textbf{专业知识深度不足}。研究表明，医疗LLM的幻觉问题并非偶发，而是LLM的一种内在理论属性，模型可能生成听起来合理但实际不准确乃至危险的医疗建议。在临床文档生成任务中，某基准研究观测到1.47\%的幻觉率和3.45\%的遗漏率，在高风险医疗场景中这类错误可能导致误诊、延误治疗乃至患者伤害。此外，通用LLM缺乏对中国医疗语境、诊疗规范和药物知识的深度理解，在专业性要求极高的慢病管理领域难以直接部署\todocite{}\todocite{}\todocite{}。

\noindent\textbf{垂类医疗大模型的微调策略与代表性成果}

针对上述局限性，医疗垂类大模型的研究路径主要包括领域数据预训练、指令微调（Instruction Fine-Tuning）和基于人类反馈的强化学习（RLHF）三类。MedChatZH是一个面向中医问答的对话模型，基于LLaMA架构的Transformer解码器并使用中医数据集进行微调，在中文医疗问答任务上取得了显著提升。讯飞星火医疗大模型以``一体两翼''技术框架（医学知识自学习为底座，认知推理与感知交互为两翼）为核心，已应用于讯飞晓医APP等产品的诊后康复管理场景。微医医疗大模型在医学知识问答、复杂医学推理等多维评测中综合得分达94.7分，位于自测榜首。2023年被业界视为中国医疗大模型发展的元年，各类垂类模型开始在临床辅助决策、健康管理等多场景中落地应用\todocite{}\todocite{}\todocite{}\todocite{}\todocite{}。

\noindent\textbf{LLM在非结构化医疗数据处理中的进展}

LLM在医疗文本的非结构化处理方面展现出独特优势，涵盖病历摘要生成、医疗实体识别（NER）、语义理解等任务。检索增强生成（RAG）技术的引入，进一步解决了LLM知识更新滞后和幻觉风险的问题------通过将LLM与外部动态知识库相结合，使模型能够实时检索最新临床指南信息，显著提升回答的准确性和可信度。一项针对糖尿病教育场景的实证研究表明，引入RAG框架（RISE）后，GPT-4、Claude 2、Google Bard三款模型的准确回答率平均提升12\%，其中Claude 2提升幅度高达19\%。面向医疗领域的MedRAG方案进一步引入知识图谱辅助推理，通过构建四层分级诊断知识图谱，在误诊率降低方面优于现有最先进的RAG方法\todocite{}\todocite{}\todocite{}。

\subsection{智能体与多智能体系统研究现状}

\noindent\textbf{智能体的核心架构}

智能体（Agent）的基本工作框架被广泛描述为``感知---推理---行动''（Perception-Reasoning-Action）的闭环循环。在医疗场景中，这一架构被进一步拓展为包含四个核心组件的体系：\textbf{感知}（处理患者输入、EHR数据等多源信息）、\textbf{规划}（分解任务目标、制定干预策略）、\textbf{行动}（执行随访、推送建议、触发预警）和\textbf{记忆}（存储历史交互与患者数据以支持个性化决策）。记忆机制对于慢病管理场景具有特殊重要性------个性化健康画像的持续构建和更新，依赖于Agent对患者长期行为数据的高效存储与检索。研究表明，将个人记忆（存储用户输入与交互历史）与系统记忆（记录中间推理状态）相结合的混合记忆架构，是实现持续个性化的关键机制\todocite{}\todocite{}\todocite{}\todocite{}。

\noindent\textbf{多智能体协同在医疗决策中的代表性工作}

\textbf{MDAgents}（Medical Decision-making Agents）是MIT、Google Research与首尔国立大学医院联合开发的自适应多Agent框架，其核心创新在于根据医疗任务的复杂度动态分配单Agent独立决策或团队协作模式，在10个医疗基准测试中的7个上超越了现有方法，精度最高提升4.2\%。该框架模拟了现实中临床医生从单科诊断到多学科会诊的分层决策过程，为复杂病例处理提供了可扩展的协作机制\todocite{}\todocite{}。

\textbf{RareAgents}是清华大学与北京协和医院联合开发、面向罕见病诊断的首个LLM驱动的多学科团队决策支持工具，其架构包含多学科团队（MDT）协调机制、动态长时记忆机制和医疗工具调用三大核心模块。具体而言，RareAgents设置了主治医师Agent和多名专科医师Agent（涵盖血液科、神经科、消化科等），通过多轮讨论机制达成诊断和治疗方案共识；动态长时记忆机制则通过嵌入历史病例相似性检索，实现个体化的记忆存储与调用。实验结果表明，RareAgents在罕见病诊断与治疗任务上优于GPT-4o和当前最先进的Agent框架，并已作为AAAI 2026口头报告被接收\todocite{}\todocite{}。

\textbf{DeepRare}（Nature Medicine，2026）则进一步采用三层架构：中央LLM主机（负责内存协调）、专业Agent服务器（负责表型与基因型分析）和外部医学资源层，通过自反式循环迭代验证假设、降低过度诊断和LLM幻觉。上述工作共同揭示了多Agent框架相比单模型的显著优势：通过专业化分工与协作，系统能够实现更准确的推理、更全面的信息整合和更可靠的决策输出\todocite{}。

\noindent\textbf{LangGraph框架}

在工程实现层面，LangGraph是当前构建多Agent应用的主流框架之一。LangGraph基于LangChain扩展，采用有向图（Graph）的方式建模Agent工作流：节点（Node）代表功能模块、边（Edge）表示信息流转方向、条件边（Conditional Edge）实现动态路由控制，状态（State）贯穿整个工作流实现信息共享与上下文传递。与Autogen、CrewAI等框架相比，LangGraph的核心优势在于其支持循环和递归（cycles and recursion）的有状态Agent工作流，天然适合慢病管理场景中需要迭代决策和长期记忆的应用需求。AWS在Amazon Bedrock上基于LangGraph构建的多Agent系统实践表明，该框架内置的状态管理与检查点机制能够保障流程连续性，条件路由支持使动态工作流调整成为可能\todocite{}\todocite{}\todocite{}\todocite{}。

\subsection{智能体评估与可解释性研究现状}

\noindent\textbf{医疗Agent的评估框架}

医疗AI系统的评估面临远比通用AI更严苛的要求，需要兼顾自动化指标与专家审核的双重标准。研究表明，Human-in-the-loop（HITL）AI在诊断成像、临床决策支持、患者监测等多个医疗应用场景中展现出相比全自动AI和传统方法的综合优势------不仅提升了诊断准确性、降低了医疗差错，还增强了临床医生对AI系统的信任。然而，HITL模式的实施需要解决电子健康档案（EHR）互操作性、责任归属框架、适应性培训协议等配套问题。临床AI安全框架的研究识别出10类失效模式，包括幻觉诊断（高严重性）、危险剂量（关键性）、紧急情况低估（关键性）等，为医疗Agent的安全评估提供了系统化分类依据。阻止有害输出抵达人工评估员之前的临床防火墙设计------包括药物剂量验证、真实患者数据模式检测、范围违规标记等------已被证明是部署医疗AI不可或缺的安全机制\todocite{}\todocite{}。

\noindent\textbf{可解释性AI与老年用户信任构建}

在面向老年患者的数字健康场景中，可解释性（Explainability）不仅是技术伦理的要求，更是推动实际使用的关键驱动力。研究表明，当AI系统无法对其推荐结果作出清晰解释时，许多临床医生和患者会产生强烈的信任焦虑，进而拒绝使用相关功能。一项面向老年人E-health交互体验的混合方法研究发现，XAI辅助的电子健康界面能够通过识别老年人偏好、提供直观可视化与简明解释，有效弥合与年龄相关的数字鸿沟。具体而言，\textbf{直观可视化}和\textbf{简明解释}是对老年用户友好的HCI（人机交互）工具设计的关键要素，而透明度和可理解性则是提升老年用户对AI驱动医疗解决方案信任度的核心路径。在医疗机构层面，将偏见透明度、模型性能可视化和主动预警等XAI策略嵌入患者端工具的医疗机构，在患者依从性和满意度指标上取得了更好的结果\todocite{}\todocite{}\todocite{}。

\noindent\textbf{``数字信任鸿沟''与适老化设计}

老年糖尿病患者在采用数字健康技术时面临独特障碍，包括技术素养不足、视力和手部灵活性下降导致的操作困难、以及对数字系统安全性的天然警惕。\emph{Diabetes Care}期刊的综述指出，老年糖尿病患者是一个在临床异质性、护理优先级和技术整合需求上与普通成人明显不同的特殊亚群，科学证据和临床实践对这一群体的数字健康工具关注仍然不足。针对这一群体，循序渐进的引导式设备设计（guided devices）被认为是帮助老年人克服技术障碍的根本解决方案，而语音交互作为最自然的人机交互形式，被认为是降低老年用户技术门槛的优先路径之一\todocite{}\todocite{}\todocite{}。

\subsection{现有糖尿病管理系统与平台述评}

\noindent\textbf{代表性系统功能分析}

目前国内已有多款面向糖尿病管理的数字化平台和移动应用上线运营，其中以智云健康和讯飞晓医最具代表性。\textbf{智云健康}（ClouDr）已获哈佛医学院Joslin糖尿病中心认证，以院内外一体化慢病管理为核心，通过智云健康APP（患者端）和智云医生APP（医生端）实现院内诊疗数据与院外居家监测数据的打通，提供互联网问诊、开方、拿药及个性化慢病管理方案等服务。根据平台公布的数据，规范使用智云健康的患者糖化血红蛋白降低效果优于国外知名慢病管理平台Livongo（0.8\%）。\textbf{讯飞星火医疗大模型}则以``医学知识自学习''为底座，构建认知推理与感知交互双翼，已在诊后康复管理、智能随访等场景中应用。适老化智能终端方面，天猫精灵颐养版、微信关怀模式等产品通过简化界面、放大字体、语音交互等设计降低老年人使用门槛，但这些终端尚未与专业糖尿病管理系统深度整合\todocite{}\todocite{}\todocite{}。

\noindent\textbf{现有系统的共同不足与本研究的切入点}

通过对现有系统的梳理分析，可以归纳出以下共同不足：\textbf{第一，AI智能化程度不足}------现有平台主要提供数据记录、推送提醒等被动式功能，缺乏基于多Agent协同的主动风险评估与预警能力；\textbf{第二，缺乏多智能体协同决策机制}------现有系统多为单一AI模型或规则引擎驱动，难以模拟真实临床多学科协作的复杂决策逻辑；\textbf{第三，适老化设计深度不足}------语音交互功能支持有限，可解释性呈现方式对认知能力下降的老年患者不够友好；\textbf{第四，知识库与指南更新机制薄弱}------多数系统的知识体系与最新临床指南（如《中国糖尿病防治指南2024版》）的同步机制不健全，缺乏RAG支撑的动态知识检索能力。

本研究的切入点正在于填补上述空缺：以LLM+多Agent+RAG为技术核心，以《中国糖尿病防治指南2024版》为知识规则基础，以适老化语音交互为接入方式，构建一个面向社区基层的、具备主动风险评估与预警能力的老年糖尿病智能随访管理系统，在现有平台的基础上实现技术能力的整体跃升。


% ------------------------------------------------------------------
\section{相关技术与理论基础}

\subsection{2型糖尿病管理的医学知识基础}

\noindent\textbf{核心监测指标体系}

2型糖尿病的临床管理依赖一套多维度监测指标体系，各指标在评估患者代谢状况和并发症风险方面发挥不可替代的作用。\textbf{空腹血糖}（Fasting Plasma Glucose，FPG）是最基础的血糖评估指标，正常范围为3.9---6.1 mmol/L，诊断糖尿病的切点为$\geq$7.0 mmol/L。\textbf{餐后2小时血糖}（2-Hour Postprandial Glucose，2hPG）反映机体餐后葡萄糖代谢能力，诊断糖尿病的切点为$\geq$11.1 mmol/L。\textbf{糖化血红蛋白}（HbA1c）是过去2---3个月平均血糖水平的综合反映，是慢病长期管理质量评估的``金标准''，《中国糖尿病防治指南（2024版）》建议大多数非妊娠成人2型糖尿病患者的HbA1c控制目标为${<}7.0\%$，但对于高龄老年、预期寿命有限或多病共存的患者，应适当放宽目标至7.5\%---8.5\%不等。此外，\textbf{BMI}、\textbf{血压}和\textbf{血脂}（LDL-C、TG、HDL-C）是评估心血管代谢综合风险的关键辅助指标\todocite{}\todocite{}。

\noindent\textbf{并发症筛查指标}

糖尿病并发症的早期筛查是社区管理的重要职能，主要包括三个维度。\textbf{肾病筛查}：通过检测尿微量白蛋白肌酐比（UACR）和估算肾小球滤过率（eGFR），每年至少进行一次系统评估；《初级医疗2型糖尿病指南（2025版）》推荐基层医疗机构每年常规进行尿UACR和血肌酐检测。\textbf{眼底筛查}：糖尿病视网膜病变是糖尿病患者失明的主要原因，基层医疗机构需每年完成初次全面眼底检查。\textbf{糖尿病足筛查}：北京大学第一医院5年随访队列研究显示，糖尿病患者的足部溃疡发生率高达25.83\%；基层随访应常规进行双足视触诊，检查足部畸形、皮肤色温变化及外周血管状况。上述筛查指标体系将作为本系统并发症风险评估Agent的核心规则库来源\todocite{}\todocite{}。

\noindent\textbf{《中国2型糖尿病防治指南（2024版）》核心规则提炼}

《中国糖尿病防治指南（2024版）》于2025年1月由中华医学会糖尿病学分会正式发布，共20章124页，是本系统知识库构建的权威依据。指南的核心管理规则可提炼为以下四个层面\todocite{}：

（1）\textbf{个体化血糖控制目标}：基于患者年龄、预期寿命、并发症状态、认知功能等因素进行分层设定。对于中等健康状态（预期寿命5---10年）的老年患者，HbA1c目标为$\leq$8.0\%；对于衰弱或多病共存（预期寿命${<}$5年）的患者，目标放宽至HbA1c$\leq$8.5\%，以预防严重高血糖及低血糖为主\todocite{}。

（2）\textbf{社区随访频率规范}：血糖达标的患者每3个月随访一次，血糖未达标的患者每2---4周随访一次，并每年进行一次全面并发症筛查\todocite{}。

（3）\textbf{用药阶梯方案}：生活方式干预（饮食+运动）是所有患者的基础管理措施；血糖仍未达标者启动以二甲双胍为核心的单药治疗；单药不足时启动联合用药（可考虑GLP-1RA、SGLT-2i、DPP-4i等新型降糖药物）；对于口服药物治疗失效的患者，进入胰岛素治疗方案\todocite{}。

（4）\textbf{转诊标准}：出现酮症酸中毒、严重低血糖、血糖控制极差（FPG${>}$16.7 mmol/L）、或出现严重并发症（急性心脑血管事件、透析前肾病等）时，需启动及时转诊至上级医院的流程\todocite{}。

\subsection{开发语言与框架}

\noindent\textbf{Python语言及其AI生态优势}

Python凭借其简洁的语法、丰富的第三方库（NumPy、Pandas、Scikit-learn等）以及与主流AI框架（LangChain、LangGraph、Hugging Face Transformers等）的天然兼容性，已成为AI医疗应用开发的首选语言。Python的AI生态系统为本系统提供了从数据处理、LLM调用、多Agent编排到向量检索的全栈支撑能力。

\noindent\textbf{Django Web框架}

Django是Python生态中成熟的企业级Web框架，内置ORM（对象关系映射）、Admin后台、CSRF安全防护机制等核心功能，提供完善的RESTful API构建支持。Django的ORM层能够屏蔽底层数据库差异，简化患者健康档案、随访记录、血糖数据等结构化信息的持久化操作；Admin后台则为系统管理员提供开箱即用的数据管理界面，降低运维成本。

\noindent\textbf{前端技术栈}

前端采用HTML5/CSS3/JavaScript技术组合，配合Bootstrap框架实现响应式布局和适老化界面设计。Bootstrap的响应式栅格系统保证系统在PC端、平板和手机端的一致呈现，其内置的字体放大、高对比度主题等无障碍设计特性有助于满足老年用户的交互需求。

\noindent\textbf{数据库选型}

系统拟采用SQLite进行开发阶段的轻量级本地存储，生产部署阶段迁移至MySQL以支持更大规模的并发访问和数据存储需求。Django的ORM层保障了数据库迁移的透明性。

\subsection{大语言模型应用技术}

\noindent\textbf{LLM选型与调用方式}

本系统在LLM选型方面综合考虑模型能力、中文支持质量、API稳定性和成本因素，主要候选模型包括：\textbf{GPT-4o}（OpenAI，强推理能力、广泛医疗场景验证）、\textbf{通义千问}（Qwen，阿里云，中文能力优秀、数据安全合规性好）和\textbf{GLM-4}（智谱AI，中文理解与生成能力强、国产化程度高）。最终选型将基于医疗问答准确率、响应延迟和API成本的实测对比结果确定，并通过统一的API调用封装层实现模型的可替换性。

\noindent\textbf{提示词工程}

提示词工程（Prompt Engineering）是控制LLM在医疗场景中行为的核心手段。\textbf{Zero-shot提示}直接向模型描述任务要求，不提供示例，适用于简单的信息查询和常规血糖解读场景；\textbf{Few-shot提示}通过在提示词中嵌入若干高质量示例，引导模型生成结构化、符合临床规范的输出，适用于并发症风险评估和转诊建议生成等需要精确格式控制的场景。本系统将采用\textbf{角色设定}（Role-Playing）策略，为不同Agent指定特定的临床角色（如``你是一名具有10年经验的社区内分泌专科医生''），以约束模型的专业立场和输出风格，同时通过\textbf{输出格式约束}（JSON Schema、Markdown结构）确保Agent输出的结构化程度满足下游处理需求\todocite{}\todocite{}\todocite{}\todocite{}。

\noindent\textbf{检索增强生成}

RAG（Retrieval-Augmented Generation）通过将LLM与外部动态知识库相结合，有效解决了LLM知识截止日期限制和领域知识深度不足的问题。本系统RAG架构的核心组件包括：（1）\textbf{糖尿病专业知识库}：以《中国糖尿病防治指南（2024版）》、《初级医疗2型糖尿病指南（2025版）》及权威循证医学文献为来源，构建结构化知识语料；（2）\textbf{文本向量化与索引}：采用向量数据库（ChromaDB或FAISS）存储知识库文本的嵌入向量（embedding），支持基于语义相似度的高效检索；（3）\textbf{查询-检索-生成链路}：当Agent接收到患者健康查询时，首先对查询进行语义向量化，检索Top-K最相关的知识片段，再将检索结果拼接至LLM提示词中，引导模型基于权威知识生成有依据的回答，从根本上降低幻觉风险\todocite{}\todocite{}\todocite{}\todocite{}。

\subsection{多智能体系统核心技术}

\noindent\textbf{智能体定义与基本架构}

在本研究的语境下，智能体（Agent）被定义为一个能够感知环境状态、自主制定决策计划并采取行动的计算实体。其核心架构为感知---推理---行动的闭环，增强版架构还包括记忆存储机制：\textbf{感知层}负责接收并解析患者输入（语音/文本）、结构化健康数据（血糖值、用药记录）等多源信息；\textbf{推理层}调用LLM完成风险评估、方案生成等认知任务，同时检索RAG知识库获取规则支撑；\textbf{行动层}执行具体操作，如生成随访建议、触发预警通知、更新健康画像；\textbf{记忆层}维护患者的长期健康档案，为个性化推理提供历史上下文\todocite{}\todocite{}\todocite{}。

\noindent\textbf{LangGraph框架架构原理}

LangGraph是本系统多Agent编排的核心技术框架。其基本构成要素包括：\textbf{图}（Graph）------定义Agent工作流的整体拓扑结构；\textbf{节点}（Node）------代表单个Agent或功能模块（如``血糖评估节点''、``转诊判断节点''）；\textbf{边}（Edge）------定义节点间的信息传递方向；\textbf{条件边}（Conditional Edge）------基于中间状态的动态分支路由，实现``如果血糖严重超标则路由至预警Agent''等条件逻辑；\textbf{状态}（State）------贯穿整个工作流的共享信息容器，确保各Agent能够访问患者的完整上下文。LangGraph支持循环（cycles）结构，允许Agent在获得新信息后重新进入推理流程，这对于需要迭代验证的复杂临床决策场景至关重要\todocite{}\todocite{}\todocite{}\todocite{}。

\noindent\textbf{Agent间通信与任务分配}

本系统规划的多Agent通信策略遵循``中央协调---专科执行''的层级架构，类似于RareAgents的主治医师---专科医师协作模式：\textbf{协调Agent}接收患者查询，判断任务复杂度并分配给对应专业Agent；\textbf{血糖评估Agent}负责血糖指标分析与HbA1c趋势判断；\textbf{并发症筛查Agent}负责肾病、眼底、足部风险的综合评估；\textbf{用药管理Agent}负责检查多药联用风险和提供用药建议；\textbf{随访预约Agent}负责根据指南规范确定下次随访时间并推送提醒。各Agent的输出将通过结构化状态传递至下一节点，最终由汇总Agent整合生成面向基层医生的综合随访报告\todocite{}\todocite{}\todocite{}。

\subsection{语音识别与自然语言处理技术}

\noindent\textbf{ASR语音识别技术}

自动语音识别（Automatic Speech Recognition，ASR）技术是本系统适老化设计的核心接入模块，旨在帮助老年患者通过口语方式完成血糖录入和健康咨询，消除键盘输入障碍。现代端到端ASR模型（如OpenAI Whisper、阿里云实时语音转写等）在中文识别上已达到较高的准确率，但研究表明，医疗领域口语中的专业实体（药物名称、检查指标等）的识别错误率仍显著高于通用文本的词错率（WER），这要求本系统在ASR之后增加医疗领域专属的后处理纠错模块\todocite{}。

\noindent\textbf{医疗实体识别与意图分类}

ASR输出的口语文本需要经过自然语言处理（NLP）管道进行结构化转化。医疗命名实体识别（Medical NER）负责从转录文本中抽取患者自述中的关键医疗实体，包括血糖数值、药物名称、症状描述、检查结果等。意图分类模块则负责判断患者话语的类别（血糖录入/用药咨询/不适症状上报/一般健康咨询等），以路由至相应的处理Agent。这一``口语$\to$结构化医疗数据''的转化链路是实现低门槛适老化交互的技术基础\todocite{}\todocite{}。

\noindent\textbf{非结构化口语到结构化医疗数据的转化链路}

完整的语音数据处理链路设计为：语音采集$\to$ASR转录$\to$医疗NER抽取$\to$意图分类$\to$实体标准化（如将``空腹血糖七点八''转化为FPG=7.8 mmol/L）$\to$写入患者健康档案。LLM在这一链路中扮演双重角色：既作为NER和意图分类的推理引擎（通过Few-shot提示实现），也作为歧义消解和语义理解的智能中枢，处理口语化、非规范表达（如``血糖还行''、``没怎么量''等）中的语义不确定性。

\subsection{面向对象开发方法}

\noindent\textbf{面向对象分析与设计基本原则}

本系统的软件工程方法论采用面向对象分析与设计（Object-Oriented Analysis and Design，OOA/OOD）范式。OOA阶段通过对社区糖尿病管理业务流程的需求建模，识别系统中的核心领域对象（如Patient、HealthRecord、VisitRecord、GlucoseReading、Agent等）及其属性与行为；OOD阶段则围绕SOLID原则（单一职责、开闭、里氏替换、接口隔离、依赖倒置）进行类设计，以确保系统的可维护性、可扩展性和可测试性。多Agent系统与面向对象设计的深度契合性体现在：每个Agent可被建模为一个封装了感知、推理和行动能力的高内聚对象，Agent间的通信对应对象间的消息传递。

\noindent\textbf{UML建模在本系统中的应用策略}

统一建模语言（UML）作为系统分析与设计的可视化语言，将在本系统中应用于以下层面：（1）\textbf{用例图}（Use Case Diagram）------描述患者、基层医生、系统管理员与系统的交互场景；（2）\textbf{类图}（Class Diagram）------建模核心领域对象及其关联关系，包括Patient、Agent、HealthProfile等核心类的属性与方法；（3）\textbf{顺序图}（Sequence Diagram）------刻画多Agent协同随访决策的消息时序流程；（4）\textbf{活动图}（Activity Diagram）------描述血糖评估、转诊判断等关键业务流程的控制逻辑。UML建模将贯穿系统设计章节，为系统的实现提供清晰的蓝图指引。


% ------------------------------------------------------------------
\section{研究内容与论文结构}

\textcolor{red}{[TODO: 补充本节，概述论文各章内容安排]}

本文共分为七章，各章内容安排如下：

第一章为绪论，阐述研究背景、研究意义、国内外研究现状及相关技术基础。

第二章为相关技术介绍，详细介绍本系统涉及的关键技术原理。\textcolor{red}{[TODO: 视是否独立成章调整]}

第三章为系统需求分析，从功能需求和非功能需求两个维度对系统进行分析，并通过用例图和用例描述明确系统的功能边界。

第四章为系统设计，包括系统总体架构设计、多智能体协同设计、数据库设计和接口设计。

第五章为系统实现，展示核心功能模块的实现过程和关键代码。

第六章为系统测试，对系统进行功能测试和性能测试，验证系统的正确性和可用性。

第七章为总结与展望，总结全文工作，分析不足之处，并对未来研究方向进行展望。
